\chapter{Statistiques \& Probabilités}
\label{chap:stats_probas}

% ==============================================================================
% 1. OBJECTIFS & COMPÉTENCES DU RÉFÉRENTIEL
% ==============================================================================
\begin{cadreretenu}[Objectifs \& Compétences du DNB Pro]
\begin{itemize}
    \item \textbf{Statistiques} : Effectifs, fréquences, moyenne arithmétique simple et pondérée, médiane, étendue, diagrammes en bâtons et circulaires.
    \item \textbf{Probabilités} : Vocabulaire des événements (certain, impossible, contraire), calcul de probabilités simples, règle de la somme $P(A) + P(\bar{A}) = 1$, arbres de probabilités à deux épreuves.
    \item \textbf{Capacités} :
    \begin{itemize}
        \item \capRealiser\ Calculer la moyenne et la médiane d'une série statistique.
        \item \capAnalyser\ Comparer deux séries statistiques à l'aide des indicateurs de position et de dispersion.
        \item \capRealiser\ Déterminer la probabilité d'un événement sous forme de fraction irréductible, décimale ou pourcentage.
        \item \capValider\ Exploiter un arbre pondéré ou un tableau à double entrée pour dénombrer des issues.
    \end{itemize}
\end{itemize}
\end{cadreretenu}

\vspace{0.3cm}

% ==============================================================================
% 2. COURS : STATISTIQUES À UNE VARIABLE
% ==============================================================================
\section{Cours : Statistiques à Une Variable}

\subsection{Indicateurs de Position : Moyenne et Médiane}

\begin{cadredef}[Moyenne pondérée]
La \textbf{moyenne} pondérée d'une série statistique est le quotient de la somme de toutes les valeurs multipliées par leurs effectifs respectifs par l'effectif total $N$ :
\[ \bar{x} = \frac{n_1 x_1 + n_2 x_2 + \dots + n_p x_p}{N} = \frac{\sum n_i x_i}{N} \]
\end{cadredef}

\begin{cadredef}[Médiane et Étendue]
\begin{itemize}
    \item La \textbf{médiane} $Me$ est une valeur qui partage la série ordonnée en deux groupes de même effectif (au moins 50\,\% des valeurs sont inférieures ou égales à $Me$).
    \item L'\textbf{étendue} $E$ est la différence entre la plus grande valeur et la plus petite valeur :
    \[ E = x_{\max} - x_{\min} \]
\end{itemize}
\end{cadredef}

\begin{cadreexemple}[Calcul complet sur une série]
Contrôle qualité de 10 pièces usinées (longueur en mm) : $48\,;\; 49\,;\; 50\,;\; 50\,;\; 50\,;\; 51\,;\; 51\,;\; 52\,;\; 53\,;\; 56$.
\begin{itemize}
    \item \textbf{Moyenne} : $\bar{x} = \frac{48 + 49 + 3\times 50 + 2\times 51 + 52 + 53 + 56}{10} = \frac{510}{10} = \repbox{51\text{ mm}}$.
    \item \textbf{Médiane} : L'effectif total est $N = 10$ (pair). $10 / 2 = 5$. La médiane est la moyenne entre la 5\textsuperscript{e} valeur ($50$) et la 6\textsuperscript{e} valeur ($51$) : $Me = \frac{50+51}{2} = \repbox{50{,}5\text{ mm}}$.
    \item \textbf{Étendue} : $E = 56 - 48 = \repbox{8\text{ mm}}$.
\end{itemize}
\end{cadreexemple}

% ==============================================================================
% 3. COURS : PROBABILITÉS
% ==============================================================================
\section{Cours : Probabilités}

\begin{cadreprop}[Propriétés fondamentales]
Pour tout événement $A$ d'un univers $\Omega$ :
\begin{itemize}
    \item $0 \le P(A) \le 1$.
    \item En situation d'\textbf{équiprobabilité} :
    \[ P(A) = \frac{\text{Nombre d'issues favorables à } A}{\text{Nombre total d'issues possibles}} \]
    \item Événement contraire $\bar{A}$ : $P(\bar{A}) = 1 - P(A)$.
\end{itemize}
\end{cadreprop}

% ==============================================================================
% 4. QUESTIONS FLASH / AUTOMATISMES
% ==============================================================================
\section{Questions Flash / Automatismes}

\begin{automatisme}[Calculs Rapides (10 min)]
\noindent
\begin{tabularx}{\linewidth}{X|X}
\textbf{1.} Moyenne de $12$, $14$ et $16$ : $\bar{x} = \pointille[1.5cm]$ & \textbf{4.} On lance un dé équilibré à 6 faces. $P(\text{obtenir un nombre pair}) = \pointille[1.5cm]$ \\
\textbf{2.} Étendue de la série $\{ 3 ; 8 ; 15 ; 2 ; 19 \}$ : $E = \pointille[1.5cm]$ & \textbf{5.} Si $P(A) = 0{,}35$, alors $P(\bar{A}) = \pointille[1.5cm]$ \\
\textbf{3.} Médiane de la série ordonnée $\{ 4 ; 7 ; 9 ; 12 ; 15 \}$ : $Me = \pointille[1.5cm]$ & \textbf{6.} Une probabilité peut-elle être égale à $1{,}2$ ? \pointille[1.5cm] \\
\end{tabularx}
\end{automatisme}

% ==============================================================================
% 5. TRAVAUX DIRIGÉS (TD)
% ==============================================================================
\section{Travaux Dirigés (TD) : Exercices d'Entraînement}

\begin{exercice}[4 pts]{\capRealiser}{Contrôle Statistique en Logistique}
Un logisticien mesure les temps de préparation de 50 commandes dans un entrepôt (en minutes) :
\begin{center}
\begin{tabular}{|l|c|c|c|c|c|}
\hline
\rowcolor{slate200} \textbf{Temps (min)} & 5 & 10 & 15 & 20 & 25 \\
\hline
\textbf{Effectif} & 8 & 14 & 16 & 8 & 4 \\
\hline
\end{tabular}
\end{center}
\begin{enumerate}
    \item Calculer le temps moyen $\bar{t}$ de préparation d'une commande.
    \item Déterminer la classe médiane et l'étendue de cette série.
    \item Quel pourcentage de commandes est préparé en 15 minutes ou moins ?
\end{enumerate}
\lignesreponse[5]
\end{exercice}

\begin{exercice}[4 pts]{\capRealiser}{Tirage Aléatoire dans un Stock de Pièces}
Dans une boîte d'un atelier mécanique, on trouve :
\begin{itemize}
    \item 15 vis à tête fraisée (dont 5 défectueuses) ;
    \item 25 vis à tête cylindrique (dont 3 défectueuses).
\end{itemize}
On tire une vis au hasard dans la boîte.
\begin{enumerate}
    \item Quelle est la probabilité que la vis soit à tête fraisée ?
    \item Quelle est la probabilité que la vis soit défectueuse ?
    \item Quelle est la probabilité que la vis soit conforme (non défectueuse) ?
\end{enumerate}
\lignesreponse[5]
\end{exercice}

% ==============================================================================
% 6. SUJET DNB PRO
% ==============================================================================
\section{Vers le Brevet (DNB Pro)}

\begin{sujetdnb}[DNB Pro Métropole][9 pts]{Consommation Électrique \& Panneaux Photovoltaïques}
Une entreprise artisanale installe des panneaux solaires sur le toit de son atelier.
Le tableau suivant donne la production électrique mensuelle mesurée sur l'année (en kWh) :

\begin{center}
\begin{tabular}{|c|c|c|c|c|c|c|c|c|c|c|c|c|}
\hline
\rowcolor{slate200} \textbf{Mois} & J & F & M & A & M & J & J & A & S & O & N & D \\
\hline
\textbf{kWh} & 180 & 240 & 420 & 580 & 750 & 820 & 860 & 810 & 630 & 410 & 220 & 160 \\
\hline
\end{tabular}
\end{center}

\begin{enumerate}
    \item \capRealiser\ Calculer la production moyenne mensuelle d'électricité sur l'année.
    \item \capRealiser\ Déterminer la médiane et l'étendue de cette production.
    \item \capAnalyser\ On choisit un mois de l'année au hasard.
    \begin{enumerate}
        \item Quelle est la probabilité que la production dépasse $600\text{ kWh}$ ?
        \item Quelle est la probabilité que la production soit inférieure à $200\text{ kWh}$ ?
    \end{enumerate}
\end{enumerate}
\lignesreponse[6]
\end{sujetdnb}

% ==============================================================================
% 7. CORRIGÉS DÉTAILLÉS
% ==============================================================================
\section{Corrigés Détaillés du Chapitre}

\begin{solution}[Corrigé des Exercices et du Sujet DNB]
\textbf{Exercice 1 (Logistique)} :
1. Total effectif $N = 8 + 14 + 16 + 8 + 4 = 50$. \\
$\bar{t} = \frac{5\times 8 + 10\times 14 + 15\times 16 + 20\times 8 + 25\times 4}{50} = \frac{40 + 140 + 240 + 160 + 100}{50} = \frac{680}{50} = \repbox{13{,}6\text{ min}}$. \\
2. $50/2 = 25$. Effectifs cumulés : $8, 22, 38$. La médiane est de \textbf{15 min}. Étendue $E = 25 - 5 = \repbox{20\text{ min}}$. \\
3. Nombre de commandes $\le 15\text{ min}$ : $8 + 14 + 16 = 38$. Soit $\frac{38}{50} = \repbox{76\,\%}$.

\vspace{2mm}
\textbf{Exercice 2 (Stock de pièces)} :
Total de vis $= 15 + 25 = 40$. \\
1. $P(\text{fraisée}) = \frac{15}{40} = \frac{3}{8} = \repbox{0{,}375}$ (soit 37,5\,\%). \\
2. Total défectueuses $= 5 + 3 = 8$. $P(\text{défectueuse}) = \frac{8}{40} = \frac{1}{5} = \repbox{0{,}20}$ (soit 20\,\%). \\
3. $P(\text{conforme}) = 1 - 0{,}20 = \repbox{0{,}80}$ (soit 80\,\%).

\vspace{2mm}
\textbf{Sujet DNB Pro (Solaire)} :
1. Total annuel $= 180 + 240 + 420 + 580 + 750 + 820 + 860 + 810 + 630 + 410 + 220 + 160 = 6\,080\text{ kWh}$. \\
Moyenne mensuelle $= 6\,080 / 12 \approx \repbox{506{,}67\text{ kWh}}$. \\
2. Série ordonnée (12 valeurs) : $160, 180, 220, 240, 410, 420, 580, 630, 750, 810, 820, 860$. \\
Médiane : moyenne entre 6\textsuperscript{e} ($420$) et 7\textsuperscript{e} ($580$) : $Me = \frac{420+580}{2} = \repbox{500\text{ kWh}}$. \\
Étendue : $E = 860 - 160 = \repbox{700\text{ kWh}}$. \\
3. (a) Mois $> 600\text{ kWh}$ : Juin, Juillet, Août, Septembre, Mai (soit 5 mois). $P = \frac{5}{12} \approx \repbox{0{,}417}$. \\
(b) Mois $< 200\text{ kWh}$ : Janvier, Décembre (soit 2 mois). $P = \frac{2}{12} = \frac{1}{6} \approx \repbox{0{,}167}$.
\end{solution}
